\chapter{ซอร์สโค้ดเฟิร์มแวร์ C++ สำหรับไมโครคอนโทรลเลอร์ ESP32}
\label{app:firmware}

รหัสต้นแบบเฟิร์มแวร์ภาษา C++ สำหรับนำไปคอมไพล์และแฟลชลงสู่บอร์ดไมโครคอนโทรลเลอร์ ESP32 หรือ ESP32-S3 เพื่อทำหน้าที่เป็นอุปกรณ์ประมวลผลสัญญาณอิสระ (Standalone IoT Edge Node) ที่ฝังค่าน้ำหนักของแบบจำลอง AI ล่าสุดไว้ในชิป

\begin{codebox}[language=C++]{เฟิร์มแวร์ ESP32 C++ พร้อมฟังก์ชันอนุมานผล AI}
// ============================================================================
// SoilpHTxAI - Embedded AI Error Compensation Engine for ESP32 / Arduino
// Project: Development of a High-Accuracy Field-Portable Soil pH Meter Prototype
// Funding: Rambhai Barni Rajabhat University Research Fund 2569
// Principal Investigators:
//   Tanapat Tirawoot (50%), Asst.Prof.Dr. Chewa Thassana (30%),
//   Assoc.Prof.Dr. Nuntaporn Moonrungsee (10%), Assoc.Prof.Dr. Nipat Piamarun (10%)
// ============================================================================

#include <Arduino.h>
#include <math.h>

// AI Model Polynomial Coefficients (Trained via On-Device OLS Regression)
static const float C_BIAS = 6.931715f;
static const float C_E    = -0.03481958f;
static const float C_T    = 0.00956416f;
static const float C_ET   = -3.47704e-05f;
static const float C_E2   = 7.53126e-06f;
static const float C_T2   = -1.62641e-05f;

// Forward Inference Function
float predictAiCompensatedPh(float potentialMv, float tempC) {
  float e = potentialMv;
  float t = tempC;
  
  float phPred = C_BIAS +
                 (C_E * e) +
                 (C_T * t) +
                 (C_ET * e * t) +
                 (C_E2 * e * e) +
                 (C_T2 * t * t);

  if (phPred < 0.0f) phPred = 0.0f;
  if (phPred > 14.0f) phPred = 14.0f;
  return phPred;
}

void setup() {
  Serial.begin(115200);
  Serial.println("SoilpHTxAI Embedded Edge AI Firmware Initialized");
}

void loop() {
  // Read analog potential from pH probe (GPIO 34) and Temperature (DS18B20)
  float rawMv = 118.5f; // Replace with analogReadMilliVolts(34)
  float tempC = 28.5f;  // Replace with sensors.getTempCByIndex(0)
  
  float phAI = predictAiCompensatedPh(rawMv, tempC);
  
  // Format JSON Output
  Serial.print("{\"potential_mv\":");
  Serial.print(rawMv);
  Serial.print(",\"temperature_c\":");
  Serial.print(tempC);
  Serial.print(",\"ai_ph\":");
  Serial.print(phAI, 2);
  Serial.println("}");
  
  delay(1000);
}
\end{codebox}

\clearpage
